% GENERATED FILE. Edit canonical lessons, then run npm run generate:books.
% canonical-source-hash: fnv1a64:609d13afc175fd22
% canonical-lessons: GU-C247-saheb, GU-C247-karmacari, GU-C247-dhandho, GU-C247-kamavu, GU-C247-mahenatu

\chapter{Earning a living}
\label{ch:gu-earning-a-living}
\hlchaptermodality{247}

\begin{chapteropening}
\textbf{By the end of this chapter:} \emph{I can talk about bosses, employees and earning money.}

Complete the last lesson of chapter 247: I can talk about bosses, employees and earning money.
\end{chapteropening}

\section[\texorpdfstring{\gu{સાહેબ} (sāheb)}{સાહેબ (sāheb)}]{\gu{સાહેબ} (sāheb) — boss, sir}

\label{lesson:GU-C247-saheb}



\begin{quote}
Before the new one: say the Gujarati for studies, then the Gujarati for to work hard.
\end{quote}

\subsection*{You'll want to know: \gu{સાહેબ}}
\textbf{\gu{સાહેબ}} — \emph{sāheb} — \textquotedblleft{}boss, sir\textquotedblright{}.

\begin{quote}
\textbf{\gu{અમારા} \gu{સાહેબ} \gu{સવારે} \gu{વહેલા} \gu{આવે} \gu{છે}.} — \emph{amārā sāheb savāre vahelā āve che} — \textquotedblleft{}Our boss comes early in the morning.\textquotedblright{}
\end{quote}

A masculine noun.

\begin{cousinweb}
\textbf{\gu{સાહેબ}} (\emph{sāheb}) is a loan from Arabic \emph{ṣāḥib}, \textquotedblleft{}companion, owner\textquotedblright{}, which reached Gujarati through Persian as a respectful word for a master. Its Arabic root is \emph{ṣ-ḥ-b} (Lane).
\end{cousinweb}

\begin{grammarlens}[title={What you've built}]
One more word for everyday life at A2.
\end{grammarlens}

\subsection*{Your turn}
\begin{itemize}
\raggedright
  \item \emph{Say it:} \emph{sāheb}
  \item \emph{Say it:} \emph{sāheb}, once more
  \item \emph{Say it:} \emph{amārā sāheb savāre vahelā āve che} — Our boss comes early in the morning.
  \item \emph{Recall:} say the Gujarati for studies, then the Gujarati for to work hard, then say \emph{sāheb} again
\end{itemize}

\begin{tcolorbox}[breakable,colback=teal!4,colframe=teal!35!black,title={Before you move on}]
What does \gu{સાહેબ} mean? (\textquotedblleft{}Boss, sir\textquotedblright{}.) Say it once more.
\end{tcolorbox}
\section[\texorpdfstring{\gu{કર્મચારી} (karmacārī)}{કર્મચારી (karmacārī)}]{\gu{કર્મચારી} (karmacārī) — employee}

\label{lesson:GU-C247-karmacari}



\begin{quote}
Before the new one: say the Gujarati for to work hard, then the Gujarati for boss.
\end{quote}

\subsection*{You'll want to know: \gu{કર્મચારી}}
\textbf{\gu{કર્મચારી}} — \emph{karmacārī} — \textquotedblleft{}employee\textquotedblright{}.

\begin{quote}
\textbf{\gu{આ} \gu{કારખાનામાં} \gu{સો} \gu{કર્મચારી} \gu{કામ} \gu{કરે} \gu{છે}.} — \emph{ā kārkhānāmā̃ so karmacārī kām kare che} — \textquotedblleft{}A hundred employees work in this factory.\textquotedblright{}
\end{quote}

A masculine noun.

\begin{cousinweb}
\textbf{\gu{કર્મચારી}} (\emph{karmacārī}) is a learned compound built from two Sanskrit words, a tatsama: \emph{karman}, \textquotedblleft{}act, work\textquotedblright{} (CDIAL 2892), and \emph{cārin}, \textquotedblleft{}moving about\textquotedblright{} (CDIAL 4765).
\end{cousinweb}

\begin{grammarlens}[title={What you've built}]
One more word for everyday life at A2.
\end{grammarlens}

\subsection*{Your turn}
\begin{itemize}
\raggedright
  \item \emph{Say it:} \emph{karmacārī}
  \item \emph{Say it:} \emph{karmacārī}, once more
  \item \emph{Say it:} \emph{ā kārkhānāmā̃ so karmacārī kām kare che} — A hundred employees work in this factory.
  \item \emph{Recall:} say the Gujarati for to work hard, then the Gujarati for boss, then say \emph{karmacārī} again
\end{itemize}

\begin{tcolorbox}[breakable,colback=teal!4,colframe=teal!35!black,title={Before you move on}]
What does \gu{કર્મચારી} mean? (\textquotedblleft{}Employee\textquotedblright{}.) Say it once more.
\end{tcolorbox}
\section[\texorpdfstring{\gu{ધંધો} (dhandho)}{ધંધો (dhandho)}]{\gu{ધંધો} (dhandho) — business, trade}

\label{lesson:GU-C247-dhandho}



\begin{quote}
Before the new one: say the Gujarati for boss, then the Gujarati for employee.
\end{quote}

\subsection*{You'll want to know: \gu{ધંધો}}
\textbf{\gu{ધંધો}} — \emph{dhandho} — \textquotedblleft{}business, trade\textquotedblright{}.

\begin{quote}
\textbf{\gu{મારા} \gu{કાકા} \gu{ચાનો} \gu{ધંધો} \gu{કરે} \gu{છે}.} — \emph{mārā kākā cāno dhandho kare che} — \textquotedblleft{}My uncle runs a tea business.\textquotedblright{}
\end{quote}

\begin{cousinweb}
The origin of \textbf{\gu{ધંધો}} (\emph{dhandho}) is not settled. CDIAL lists a Gujarati \emph{dhandho} only in the sense \textquotedblleft{}fuss\textquotedblright{}, under Sanskrit \emph{dhandha}, \textquotedblleft{}awkwardness\textquotedblright{} (CDIAL 6727), and EtymDB links the word to a Prakrit \emph{dhaṁdha}.
\end{cousinweb}

\begin{grammarlens}[title={What you've built}]
One more word for everyday life at A2.
\end{grammarlens}

\subsection*{Your turn}
\begin{itemize}
\raggedright
  \item \emph{Say it:} \emph{dhandho}
  \item \emph{Say it:} \emph{dhandho}, once more
  \item \emph{Say it:} \emph{mārā kākā cāno dhandho kare che} — My uncle runs a tea business.
  \item \emph{Recall:} say the Gujarati for boss, then the Gujarati for employee, then say \emph{dhandho} again
\end{itemize}

\begin{tcolorbox}[breakable,colback=teal!4,colframe=teal!35!black,title={Before you move on}]
What does \gu{ધંધો} mean? (\textquotedblleft{}Business, trade\textquotedblright{}.) Say it once more.
\end{tcolorbox}
\section[\texorpdfstring{\gu{કમાવું} (kamāvũ)}{કમાવું (kamāvũ)}]{\gu{કમાવું} (kamāvũ) — to earn}

\label{lesson:GU-C247-kamavu}



\begin{quote}
Before the new one: say the Gujarati for employee, then the Gujarati for business.
\end{quote}

\subsection*{You'll want to know: \gu{કમાવું}}
\textbf{\gu{કમાવું}} — \emph{kamāvũ} — \textquotedblleft{}to earn\textquotedblright{}.

\begin{quote}
\textbf{\gu{હવે} \gu{તે} \gu{સારું} \gu{કમાય} \gu{છે}.} — \emph{have te sārũ kamāy che} — \textquotedblleft{}Now he earns well.\textquotedblright{}
\end{quote}

\begin{cousinweb}
\textbf{\gu{કમાવું}} (\emph{kamāvũ}) is a tadbhava. It comes from \emph{*karmāpayati}, \textquotedblleft{}works, earns\textquotedblright{}, a reconstructed form built on Sanskrit \emph{karman}, \textquotedblleft{}work\textquotedblright{}; CDIAL glosses the Gujarati verb \textquotedblleft{}to earn\textquotedblright{} (CDIAL 2897).
\end{cousinweb}

\begin{grammarlens}[title={What you've built}]
One more word for everyday life at A2.
\end{grammarlens}

\subsection*{Your turn}
\begin{itemize}
\raggedright
  \item \emph{Say it:} \emph{kamāvũ}
  \item \emph{Say it:} \emph{kamāvũ}, once more
  \item \emph{Say it:} \emph{have te sārũ kamāy che} — Now he earns well.
  \item \emph{Recall:} say the Gujarati for employee, then the Gujarati for business, then say \emph{kamāvũ} again
\end{itemize}

\begin{tcolorbox}[breakable,colback=teal!4,colframe=teal!35!black,title={Before you move on}]
What does \gu{કમાવું} mean? (\textquotedblleft{}To earn\textquotedblright{}.) Say it once more.
\end{tcolorbox}
\section[\texorpdfstring{\gu{મહેનતુ} (mahenatu)}{મહેનતુ (mahenatu)}]{\gu{મહેનતુ} (mahenatu) — hardworking}

\label{lesson:GU-C247-mahenatu}



\begin{quote}
Before the new one: say the Gujarati for business, then the Gujarati for to earn.
\end{quote}

\subsection*{You'll want to know: \gu{મહેનતુ}}
\textbf{\gu{મહેનતુ}} — \emph{mahenatu} — \textquotedblleft{}hardworking\textquotedblright{}.

\begin{quote}
\textbf{\gu{આ} \gu{કર્મચારી} \gu{મહેનતુ} \gu{છે}.} — \emph{ā karmacārī mahenatu che} — \textquotedblleft{}This employee is hardworking.\textquotedblright{}
\end{quote}

\begin{cousinweb}
\textbf{\gu{મહેનતુ}} (\emph{mahenatu}) is built on \emph{mahenat}, \textquotedblleft{}hard work\textquotedblright{}, with the adjective suffix \emph{-u}. \emph{Mahenat} is a loan from Arabic \emph{miḥnat}, \textquotedblleft{}trial, affliction\textquotedblright{}, through Persian (Lane).
\end{cousinweb}

\begin{grammarlens}[title={What you've built}]
One more word for everyday life at A2.
\end{grammarlens}

\subsection*{Your turn}
\begin{itemize}
\raggedright
  \item \emph{Say it:} \emph{mahenatu}
  \item \emph{Say it:} \emph{mahenatu}, once more
  \item \emph{Say it:} \emph{ā karmacārī mahenatu che} — This employee is hardworking.
  \item \emph{Recall:} say the Gujarati for business, then the Gujarati for to earn, then say \emph{mahenatu} again
\end{itemize}

\begin{tcolorbox}[breakable,colback=teal!4,colframe=teal!35!black,title={Before you move on}]
What does \gu{મહેનતુ} mean? (\textquotedblleft{}Hardworking\textquotedblright{}.) Say it once more.
\end{tcolorbox}
